\documentclass{article}
\usepackage{amsmath, amssymb, amsthm}
\title{Option 4: Supremum of Union Bound}
\author{Brooke Freeman}
\date{Nov. 2026}

\begin{document}

\maketitle

\section{Option 4}
Suppose set $A$ and $B$ both have supremums, $\alpha = \sup(A)$ and $\beta = \sup(B)$.\\
Prove that $\sup(A \cup B)$ exists, and $\sup(A \cup B) \geq \max\{\alpha,\beta\}$.\\
\indent a. Can this inequality be strict? 

\section{Rewrite}
Suppose set $A$ and $B$ both have supremums, $\alpha = \sup(A)$ and $\beta = \sup(B)$.\\
Let set $X = (A \cup B) = \{x \mid x \in A \lor x \in B\}$\\
Prove $\sup(X)$ exists and $\sup(X) \geq \max\{\alpha,\beta\}$.\\
\indent a. Can this inequality be strict?

\section{Proving $\sup(X)$ exists}

\begin{proof} 
Assume $A$ is empty. Note: we know $\sup(A) = \alpha$ exists.\\
Let $\varepsilon > 0$, if $A$ is empty then by defn. of an upper bound $\alpha - \varepsilon$ is an upper bound of $A$, so\\
\indent $\alpha \leq \alpha - \varepsilon$\\
\indent $\alpha - \alpha \leq -\varepsilon$\\
\indent $0 \leq -\varepsilon$\\
\indent $0 \geq \varepsilon$\\
This contradicts $\varepsilon > 0$ and cannot be true, by the second criteria of supremum, which states $\alpha \leq$ all upper bounds, therefore $A$ is not empty.\\ (Same can be said for set $B$)\\
Since both sets, $A$ and $B$ are not empty, then by defn. of Union $X = (A \cup B)$ is also not empty, guaranteed to have a minimum of one element.\\
(Note: that only one set of a union needs not be empty for the union to not be empty).\\
\\
(Remember: \\
\indent $X = (A \cup B) = \{x \mid x \in A \lor x \in B\}$\\
\indent $\sup(A) = \alpha$ and $\sup(B) = \beta$ )\\
Let $\zeta = \max\{\alpha,\beta\}$ then, for all $x \in X$ by defn. of supremum criteria (i)\\
\indent if $x \in A, x \leq \alpha \leq \zeta$\\
\indent if $x \in B, x \leq \beta \leq \zeta$\\
then $x \leq \zeta$ always\\
Thus by defn. of upper bound, $\zeta$ is an upper bound of $X$, therefore bounded above.
\\
In conclusion by axiom of completeness, since $X$ is non-empty, and bounded above then $X$ has a least upper bound.
\end{proof}

\section{Prove $\sup(X) \geq \max\{\alpha,\beta\}$}

\begin{proof}
Let $\zeta = \max\{\alpha,\beta\}$, $\sup(X) = s$\\
Since $\zeta$ is either $\alpha$ or $\beta$:\\
If $\zeta = \alpha$, for all $a \in A$ since $A \subseteq X$ then $a \in X$ and cause $s$ upper bound of $X$ by supremum criteria (i) $a \leq s$, then $s$ is an upper bound of $A$, then by supremum criteria (ii) $\zeta = \alpha \leq s$\\
If $\zeta = \beta$, for all $b \in B$ since $B \subseteq X$ then $b \in X$ and cause $s$ upper bound of $X$ by supremum criteria (i) $b \leq s$, then $s$ is an upper bound of $B$, then by supremum criteria (ii) $\zeta = \beta \leq s$\\
Thus in either case $\zeta \leq s$
\end{proof}
In previous section, already proved $\zeta$ is an upper bound of $X$, thus by supremum criteria (ii): $s \leq \zeta$, since both directions have been proven we actually have equality. $\sup(X) = \max\{\alpha,\beta\}$


\section{Can $\sup(X) > \max\{\alpha,\beta\}$?}
No, proved equality, $\sup(X) = \max\{\alpha,\beta\}$ so can't be strict. 

\end{document}